\chapter{Mehrfaches Bilden von Schrankenmengen}
\label{b13guide:ch:iteration}

\section{Schrankenidentitäten ohne Existenzannahmen}

\begin{theorem}[Wechselseitige Schrankenbedingung]
\label{b13guide:thm:polarity}
Für $S,T\subseteq P$ gilt
\[
 S\subseteq\Lo(T)\quad\Longleftrightarrow\quad T\subseteq\Hi(S).
\]
Insbesondere ist jede Menge in beiden ihrer doppelten
Schrankenmengen enthalten:
\[
 S\subseteq\Lo(\Hi(S)),\qquad S\subseteq\Hi(\Lo(S)).
\]
\end{theorem}

\begin{theorem}[Dreifache Schrankenbildung]
\label{b13guide:thm:triple}
Für jede Teilmenge $S\subseteq P$ gilt
\[
 \Lo(\Hi(\Lo(S)))=\Lo(S),\qquad
 \Hi(\Lo(\Hi(S)))=\Hi(S).
\]
Diese Identitäten gelten auch bei leeren Schrankenmengen.
\end{theorem}

\begin{theorem}[Doppelte Schrankenbildung]
\label{b13guide:thm:double}
Für $S\subseteq T\subseteq P$ gilt
\[
 \Lo(\Hi(S))\subseteq\Lo(\Hi(T)),\qquad
 \Hi(\Lo(S))\subseteq\Hi(\Lo(T)).
\]
Für jedes $S\subseteq P$ gilt außerdem
\[
\begin{aligned}
 \Lo(\Hi(\Lo(\Hi(S))))&=\Lo(\Hi(S)),\\
 \Hi(\Lo(\Hi(\Lo(S))))&=\Hi(\Lo(S)).
\end{aligned}
\]
Die doppelte Anwendung verändert sich also bei Wiederholung nicht.
\end{theorem}

\begin{theorem}[Größte Erweiterungen mit denselben Schranken]
\label{b13guide:thm:largest-extension}
Für $S\subseteq T\subseteq P$ gilt
\[
\begin{aligned}
 \Hi(T)=\Hi(S)&\ \Longleftrightarrow\ T\subseteq\Lo(\Hi(S)),\\
 \Lo(T)=\Lo(S)&\ \Longleftrightarrow\ T\subseteq\Hi(\Lo(S)).
\end{aligned}
\]
Somit ist $\Lo(\Hi(S))$ die bezüglich Inklusion größte Erweiterung
von $S$ mit denselben oberen Schranken; $\Hi(\Lo(S))$ ist die
größte Erweiterung mit denselben unteren Schranken.
Für beliebige $S,T\subseteq P$ gilt darüber hinaus
\[
\begin{aligned}
 \Hi(S)=\Hi(T)&\ \Longleftrightarrow\ \Lo(\Hi(S))=\Lo(\Hi(T)),\\
 \Lo(S)=\Lo(T)&\ \Longleftrightarrow\ \Hi(\Lo(S))=\Hi(\Lo(T)).
\end{aligned}
\]
\end{theorem}


\begin{theorem}[Erkennen von Schrankenmengen]
\label{b13guide:thm:recognize-bounds}
Für $M\subseteq P$ sind äquivalent:
\[
 \exists S\subseteq P\;(M=\Lo(S)),\qquad M=\Lo(\Hi(M)).
\]
Ebenso sind äquivalent:
\[
 \exists S\subseteq P\;(M=\Hi(S)),\qquad M=\Hi(\Lo(M)).
\]
Die erste Gleichheit liefert im Existenzfall die Darstellung
$M=\Lo(S)$ mit $S=\Hi(M)$, die zweite die Darstellung
$M=\Hi(S)$ mit $S=\Lo(M)$.
\end{theorem}

\section{Zusammenhang mit vorhandenen Infima und Suprema}

\begin{theorem}[Doppelte Schranken als Charakterisierung]
\label{b13guide:thm:double-principal}
Für $S\subseteq P$ und $m,s\in P$ gilt
\[
\begin{aligned}
 s=\sup S&\ \Longleftrightarrow\ \Lo(\Hi(S))=\mathord\downarrow s,\\
 m=\inf S&\ \Longleftrightarrow\ \Hi(\Lo(S))=\mathord\uparrow m.
\end{aligned}
\]
\end{theorem}

\begin{theorem}[Infimum der oberen und Supremum der unteren Schranken]
\label{b13guide:thm:bounds-extrema}
Für jede Teilmenge $S\subseteq P$ gilt
\[
\begin{aligned}
 \sup S\text{ existiert}&\ \Longleftrightarrow\
       \inf\Hi(S)\text{ existiert},\\
 \inf S\text{ existiert}&\ \Longleftrightarrow\
       \sup\Lo(S)\text{ existiert}.
\end{aligned}
\]
Im jeweiligen Existenzfall gilt
\[
 \sup S=\inf\Hi(S),\qquad \inf S=\sup\Lo(S).
\]
Dies schließt $S=\varnothing$ ein.
\end{theorem}

\begin{theorem}[Vorhandene Grenzen innerhalb von Schrankenmengen]
\label{b13guide:thm:bounds-closed}
Sei $X\subseteq\Lo(S)$. Existiert $\sup X$ in $P$, so liegt
$\sup X$ in $\Lo(S)$.
Sei $Y\subseteq\Hi(S)$. Existiert $\inf Y$ in $P$, so liegt
$\inf Y$ in $\Hi(S)$.
Beide Aussagen gelten auch für leeres $X$ beziehungsweise $Y$,
sofern der genannte Wert existiert.
\end{theorem}
